\documentclass[11pt,a4paper]{article}

\usepackage[margin=1in]{geometry}
\usepackage[T1]{fontenc}
\usepackage{lmodern}
\usepackage{hyperref}
\usepackage{longtable}

\title{Sudoku Neural-to-Fundamental-Function Research\\
Research Development Log}
\author{Kanwal Rai}
\date{\today}

\begin{document}

\maketitle

\section{Purpose}

This document records the implementation history, methodological decisions,
experimental results, reproducibility information, and research conclusions
for the Sudoku neural-to-fundamental-function research project.

\section{Research Principle}

Sudoku is used as a controlled combinatorial reasoning laboratory. The
research does not claim to resolve the P versus NP problem. Instead, it
investigates whether progressively sophisticated neural architectures can
learn and expose compact structural representations of Sudoku constraints
and whether such learned structure can ultimately be converted into an
independent non-neural function or solver.

\section{Development Status}

\begin{longtable}{p{0.18\textwidth}p{0.72\textwidth}}
\textbf{Sprint} & \textbf{Status} \\
\hline
Sprint 00 & Development environment and repository bootstrap -- COMPLETE \\
Sprint 01 & Mathematical representation, board, constraints, validator -- pending \\
Sprint 02 & Exact Sudoku solver / ground-truth oracle -- pending \\
Sprint 03 & Synthetic puzzle generator and uniqueness verification -- pending \\
Sprint 04 & Dataset engineering and leakage prevention -- pending \\
Sprint 05 & MLP baseline -- pending \\
Sprint 06 & MLP generalization -- pending \\
Sprint 07 & MLP reverse engineering -- pending \\
Sprint 08 & CNN -- pending \\
Sprint 09 & RNN/LSTM/GRU -- pending \\
Sprint 10 & GNN constraint graph -- pending \\
Sprint 11 & Transformer -- pending \\
Sprint 12 & GPT-style sequence model -- pending \\
Sprint 13 & Cross-architecture representation analysis -- pending \\
Sprint 14 & Symbolic regression / equation and rule discovery -- pending \\
Sprint 15 & Independent extracted function / solver -- pending \\
Sprint 16 & Independent validation on fresh unseen distributions -- pending \\
\end{longtable}

\section{Sprint 00}

The development environment was verified under WSL Ubuntu on the Windows
D: drive. The intended repository root is:

\begin{verbatim}
/mnt/d/0001_0001/sudoku-research
\end{verbatim}

The canonical Git branch is \texttt{main}. GitHub authentication is configured
for the \texttt{raikanwalrai} account.

\subsection{Sprint 00 Completion Record}

Sprint 00 established and verified the reproducible development foundation
for the project.

\begin{itemize}
    \item Development environment: WSL Ubuntu on the Windows D: drive.
    \item Canonical local repository:
    \texttt{/mnt/d/0001\_0001/sudoku-research}.
    \item Git default branch: \texttt{main}.
    \item Git author: Kanwal Rai.
    \item GitHub account: \texttt{raikanwalrai}.
    \item Git LFS: installed and operational.
    \item GitHub repository: \texttt{raikanwalrai/sudoku-research}.
    \item Repository visibility: public.
    \item Initial repository commit:
    \texttt{9658f4c}, ``chore: bootstrap Sudoku research repository''.
    \item The initial commit was pushed successfully to
    \texttt{origin/main}.
    \item The local working tree was verified clean after the initial push.
\end{itemize}

The repository structure, research-control documents, Git ignore boundary,
local commit, remote configuration, and GitHub commit state were independently
checked before proceeding to the next sprint.

\end{document}
